\chapter{Hepatitis Panel}

\section*{What Is This Test?}
A hepatitis panel is a combination of serologic tests that detects and differentiates acute and chronic viral hepatitis, most commonly hepatitis A virus (HAV), hepatitis B virus (HBV), and hepatitis C virus (HCV). HAV is a picornavirus transmitted via the fecal-oral route causing acute self-limited hepatitis; it does not cause chronic infection. HBV is a hepadnavirus transmitted parenterally and sexually that causes acute and chronic hepatitis; chronic HBV affects approximately 296 million people worldwide. HCV is a flavivirus transmitted primarily parenterally; 75--85\% of infections become chronic, affecting approximately 58 million people globally. The acute hepatitis panel typically includes: HAV IgM, hepatitis B surface antigen (HBsAg), hepatitis B core antibody IgM (HBcAb IgM), and hepatitis C antibody (anti-HCV). The chronic hepatitis evaluation panel includes HBsAg, HBsAb (anti-HBs), HBcAb total (IgG + IgM), HBeAg, HBeAb (anti-HBe), HBV DNA, anti-HCV, and HCV RNA. Additional tests like HAV IgG, HEV IgM/IgG (hepatitis E virus, important in pregnancy and immunocompromised), and HDV serology (hepatitis D, which requires HBsAg for replication) are ordered based on clinical context.

\section*{Alternative Names}
Viral Hepatitis Serology; Hepatitis A/B/C Panel; Hep Panel; Acute Hepatitis Panel; Chronic Hepatitis Panel; HAV/HBV/HCV Screen; Hepatitis Serology; Liver Panel with Hepatitis Serology

\section*{Principle}
All hepatitis serologic markers are detected by chemiluminescent immunoassay (CLIA) or enzyme immunoassay (EIA). HAV IgM is the marker of acute HAV infection; detectable at symptom onset and remains positive for 3--6 months. HAV IgG appears during convalescence and persists for life, conferring immunity. HBsAg is the serologic hallmark of HBV infection; detectable 1--10 weeks after exposure and before symptoms. Persistence beyond 6 months defines chronic infection. HBcAb IgM is the marker of acute or recent HBV infection (window between HBsAg clearance and appearance of anti-HBs). HBcAb total (IgG + IgM) remains positive for life in anyone ever infected; it distinguishes infection-acquired from vaccine-acquired immunity (vaccine induces anti-HBs only, not anti-HBc). HBeAg is a marker of active HBV replication and high infectivity. Anti-HCV is detected by a 3rd-generation EIA or CLIA with sensitivity >97\%. All positive anti-HCV results should be confirmed with HCV RNA by PCR, as 15--25\% of anti-HCV-positive individuals spontaneously clear the virus and will have negative HCV RNA. HBV DNA and HCV RNA are quantitated by RT-PCR and reported in IU/mL. HCV genotype is determined by sequencing or line probe assay to guide choice and duration of direct-acting antiviral (DAA) therapy.

\section*{History}
Hepatitis A was distinguished from hepatitis B by Saul Krugman's human experiments at Willowbrook State School in the 1960s. HAV was visualized by electron microscopy in 1973 and the HAV vaccine was licensed in 1995. HBV was identified by discovery of the Australia antigen (now HBsAg) by Baruch Blumberg in 1965 (Nobel Prize 1976). The HBV vaccine (plasma-derived) was licensed in 1981; recombinant vaccine in 1986. HCV was discovered in 1989 by Chiron scientists, the first pathogen identified purely by molecular cloning without visualization or culture (Nobel Prize 2020 for Harvey Alter, Michael Houghton, and Charles Rice). The first HCV antibody test was licensed in 1990, enabling blood supply screening that virtually eliminated transfusion-associated HCV. Interferon-based HCV therapy (1990s--2011) was associated with cure rates of 40--50\%. The first DAAs (telaprevir, boceprevir) were introduced in 2011, and sofosbuvir in 2013 launched the all-oral DAA era with cure rates >95\%. WHO set targets for global elimination of viral hepatitis as a public health threat by 2030.

\section*{Why This Test Is Ordered}
\begin{itemize}
  \item Evaluation of patients presenting with jaundice, right upper quadrant pain, nausea/vomiting, fatigue, dark urine, clay-colored stools, or elevated serum transaminases (ALT and AST) to determine the viral etiology of acute hepatitis
  \item Screening for chronic HBV and HCV in at-risk populations: persons born in countries with HBV prevalence >2\%, injection drug users, MSM, individuals with HIV, hemodialysis patients, recipients of blood products or organ transplants before 1992, persons with elevated ALT of unknown cause, and persons born between 1945--1965 (HCV birth-cohort screening)
  \item Routine prenatal screening for HBsAg in all pregnant women (standard of care in the first trimester to prevent perinatal HBV transmission)
  \item Pre-vaccination screening for HAV and HBV immunity in at-risk individuals; post-vaccination anti-HBs titer 1--2 months after the last vaccine dose in immunocompromised individuals, healthcare workers, and hemodialysis patients to confirm vaccine response (>10 mIU/mL = protective)
  \item Baseline evaluation and ongoing monitoring of patients with chronic HBV or HCV infection, including assessment of disease activity (HBV DNA, HBeAg, ALT) or need for and response to antiviral therapy (HCV RNA before, during, and 12 weeks after DAA therapy for sustained virologic response/SVR12)
  \item Workup of patients with fulminant hepatic failure, where viral etiologies (HAV, HBV, HEV, HSV, VZV) must be rapidly excluded
  \item Organ donor and recipient screening to prevent donor-derived HBV and HCV transmission
\end{itemize}

\section*{Is This a Primary or Secondary Test}
Hepatitis serology is the primary screening test for viral hepatitis. The acute hepatitis panel (HAV IgM, HBsAg, HBcAb IgM, anti-HCV $\pm$ HCV RNA) provides rapid etiologic diagnosis. Screening for HBV and HCV in at-risk populations is primary. Confirmatory nucleic acid testing (HBV DNA, HCV RNA) is required to confirm active viral replication.

\section*{How This Test Is Performed}
Venous blood collected in an SST (serum separator) or EDTA (plasma) tube. Approximately 5--7 mL of whole blood yields sufficient serum or plasma. Serum should be separated within 2 hours of collection. Specimens can be refrigerated at 2--8\textdegree{}C for up to 7 days or frozen at -20\textdegree{}C to -80\textdegree{}C for extended storage. The initial serologic panel is typically performed on the same blood sample. Additional testing for HBV DNA or HCV RNA can be performed on the same specimen if sufficient volume is available, or on a separate EDTA plasma specimen (EDTA is preferred for viral load assays). For the comprehensive hepatitis panel, it is efficient to draw one SST tube for serology and one EDTA tube for potential viral load testing.

\section*{What Preparation Is Needed}
No fasting is required for hepatitis serology. Patients should report any prior hepatitis vaccination (HAV, HBV) and known history of hepatitis. Certain medications (high-dose biotin/supplements >5 mg/day, monoclonal antibodies) may interfere with immunoassays and should be noted; biotin supplements should be withheld for 48--72 hours before testing.

\section*{Contraindications and Precautions}
A positive HAV IgM in the absence of clinical hepatitis may represent a false positive or prolonged IgM positivity (HAV IgM can persist for 3--6 months and occasionally >1 year). The presence of rheumatoid factor or acute EBV/CMV infection may cause false-positive IgM assays. A positive anti-HCV in a low-risk individual (e.g., blood donor with no risk factors) has a false-positive rate of 15--35\%; HCV RNA testing is mandatory for confirmation. In acute HCV infection, anti-HCV may be negative for the first 6--8 weeks (window period); HCV RNA should be tested directly when acute HCV is suspected. In immunosuppressed patients, anti-HCV may be falsely negative due to impaired antibody production; HCV RNA testing is indicated if suspicion remains despite negative serology, particularly in hemodialysis patients and those with hematologic malignancies. HBsAg mutants (escape mutants with mutations in the "a" determinant of the S gene) may not be detected by some older HBsAg assays; most modern assays use capture antibodies that recognize multiple epitopes and detect the vast majority of mutants.

\section*{Risks}
Standard venipuncture risks. No direct risks from the serologic tests. The major clinical risk is the failure to diagnose chronic HBV or HCV, as both can progress silently to cirrhosis, end-stage liver disease, and hepatocellular carcinoma. Birth-cohort HCV screening and risk-based HBV screening seek to address this subclinical burden.

\section*{What Happens After the Test}
Serology results available in 1--24 hours. HBV DNA and HCV RNA results available in 24--72 hours. Positive HAV IgM, diagnostic of acute hepatitis A, requires public health notification. Positive HBsAg requires linkage to care for chronic HBV evaluation and management. Positive anti-HCV requires reflex HCV RNA testing; if HCV RNA is positive, linkage to care for DAA treatment. All cases of acute and chronic viral hepatitis are notifiable to public health authorities.

\section*{Understanding Results}

\subsection*{Below Normal}
\begin{itemize}
  \item \textbf{All hepatitis panel markers negative:} No evidence of acute or past HAV, HBV, or HCV infection. HAV and HBV IgG negative indicates susceptibility to HAV and HBV. The patient should be offered vaccination against HAV and HBV if at risk
  \item \textbf{Isolated HAV IgG positive, all other markers negative:} Past (resolved) HAV infection or HAV vaccination indicating immunity; no acute or chronic hepatitis
  \item \textbf{Anti-HBs positive (>10 mIU/mL), all other markers negative:} Immunity to HBV through vaccination. No past or current HBV infection (HBcAb is negative in vaccine-immune)
  \item \textbf{Anti-HCV negative:} No HCV infection unless (a) recent exposure within the window period (tested <6--8 weeks from exposure), or (b) immunocompromised patient with impaired antibody production. If acute HCV is clinically suspected, HCV RNA should be tested directly
  \item \textbf{HCV RNA not detected:} In an anti-HCV positive individual, indicates past (resolved) HCV infection with spontaneous clearance (15--25\% of cases) or successful treatment (sustained virologic response, SVR). In an anti-HCV negative individual, essentially excludes active HCV
\end{itemize}

\subsection*{Above Normal}
\begin{itemize}
  \item \textbf{HAV IgM positive:} Acute hepatitis A infection. Peaks during symptomatic phase and persists 3--6 months. Supportive care is the mainstay; hospitalization indicated for coagulopathy, encephalopathy, or severe nausea preventing oral intake. HAV vaccine or immune globulin should be administered to unvaccinated close contacts
  \item \textbf{HBsAg positive:} HBV infection. If concurrent with HBcAb IgM positive and clinical hepatitis, indicates acute HBV infection. If HBsAg is positive for >6 months, indicates chronic HBV. Further testing: HBeAg (marker of active replication and high infectivity), ALT (marker of hepatic inflammation), HBV DNA (quantitative viral load), and liver fibrosis assessment (FibroScan, Fib-4, liver biopsy) are indicated. Chronic HBV requires lifelong monitoring and consideration of antiviral therapy (tenofovir, entecavir)
  \item \textbf{HBsAg negative, HBcAb total positive, anti-HBs positive:} Resolved (past) HBV infection with natural immunity. No active infection. However, HBV covalently closed circular DNA (cccDNA) persists in the liver, and reactivation can occur with immunosuppression (rituximab, high-dose chemotherapy, TNF inhibitors, prolonged corticosteroids). Anti-HBV prophylaxis should be administered to HBsAg-negative/HBcAb-positive patients receiving these therapies
  \item \textbf{Anti-HCV positive, HCV RNA detected:} Active HCV infection. The HCV RNA level and HCV genotype guide DAA regimen selection. All patients with active HCV should be treated with oral DAAs (glecaprevir/pibrentasvir, sofosbuvir/velpatasvir) with cure rates >95\%. Post-treatment SVR12 (HCV RNA negative 12 weeks after therapy completion) defines cure
  \item \textbf{Anti-HCV positive, HCV RNA not detected, repeat HCV RNA negative:} Past (resolved) HCV (spontaneous clearance or cured). No current infection. A single negative HCV RNA in the setting of anti-HCV positivity should be confirmed by a second negative HCV RNA before concluding spontaneous clearance
\end{itemize}

\section*{Normal Results}
HAV: \textbf{HAV IgM Negative, HAV IgG Negative} (susceptible; candidate for vaccination) OR \textbf{HAV IgG Positive} (immune). HBV: \textbf{HBsAg Negative, HBcAb Negative, anti-HBs Negative} (susceptible) OR \textbf{anti-HBs >10 mIU/mL} (immune through vaccination or resolved infection). HCV: \textbf{Anti-HCV Negative}. If anti-HCV positive: \textbf{HCV RNA Not Detected} (resolved or false positive). Quantitative HBV DNA and HCV RNA: \textbf{Not detected} or \textbf{<LLOQ}.

\section*{Follow-up Tests}
\begin{itemize}
  \item \textbf{HCV RNA quantitative (RT-PCR) and HCV genotype:} Mandatory for any positive anti-HCV antibody result. Genotype (1a, 1b, 2, 3, 4, 5, 6) determines the choice and duration of DAA therapy. HCV RNA is monitored at week 4 of treatment, end of treatment, and 12 weeks after treatment (SVR12)
  \item \textbf{HBV DNA quantitative:} For all HBsAg-positive patients. Levels >2,000 IU/mL in HBeAg-negative chronic hepatitis, or any level combined with elevated ALT and/or significant liver fibrosis, are thresholds for initiating antiviral therapy
  \item \textbf{HBeAg and anti-HBe:} For all HBsAg-positive patients. HBeAg positivity indicates high viral replication and infectivity. Seroconversion (HBeAg loss, appearance of anti-HBe) is a favorable therapeutic endpoint in HBeAg-positive chronic hepatitis B
  \item \textbf{Liver panel (ALT, AST, alkaline phosphatase, bilirubin, albumin, INR):} For all patients with viral hepatitis to assess hepatic synthetic function and degree of liver injury. ALT >2x upper limit of normal is a marker of active liver inflammation
  \item \textbf{Abdominal ultrasound with Doppler:} For all HBsAg- or HCV RNA-positive patients to assess for cirrhosis, portal hypertension, and hepatocellular carcinoma (HCC) screening. HCC screening with ultrasound $\pm$ alpha-fetoprotein every 6 months is standard for all patients with cirrhosis
  \item \textbf{HEV serology (HEV IgM/IgG):} When acute hepatitis is negative for HAV, HBV, and HCV, particularly if the patient has traveled to HEV-endemic regions, is pregnant (high HEV mortality), or is immunocompromised (risk of chronic HEV)
  \item \textbf{HDV serology (anti-HDV):} For all HBsAg-positive patients, as hepatitis D requires HBV for replication and carries the worst prognosis of all viral hepatitides
\end{itemize}

\section*{References}
\begin{enumerate}
  \item Terrault NA, Lok ASF, McMahon BJ, et al. Update on prevention, diagnosis, and treatment of chronic hepatitis B: AASLD 2018 hepatitis B guidance. Hepatology. 2018;67(4):1560-1599.
  \item AASLD-IDSA. HCV Guidance: Recommendations for Testing, Managing, and Treating Hepatitis C. Updated periodically at www.hcvguidelines.org.
  \item CDC. Recommendations for the identification of chronic hepatitis C virus infection among persons born during 1945--1965. MMWR Recomm Rep. 2012;61(RR-4):1-32.
  \item Schillie S, Vellozzi C, Reingold A, et al. Prevention of hepatitis B virus infection in the United States: ACIP recommendations. MMWR Recomm Rep. 2018;67(1):1-31.
  \item Nelson NP, Link-Gelles R, Hofmeister MG, et al. Update: ACIP recommendations for use of hepatitis A vaccine for PEP and for PEP and for international travel. MMWR. 2019;68(42):947-950.
\end{enumerate}

\clearpage
